\chapter{Gateway（网关核心）}

\section{概述}

Gateway 网关是 OpenClaw 的核心进程，承担以下职责：

\begin{itemize}
  \item \textbf{消息路由}：接收各渠道入站消息，路由到正确的 Agent
  \item \textbf{会话管理}：维护所有会话状态（唯一事实来源）
  \item \textbf{智能体运行时}：内嵌 pi-agent-core，执行智能体循环
  \item \textbf{RPC 服务}：对外提供 WebSocket/HTTP API
  \item \textbf{插件宿主}：加载并运行所有已注册插件
\end{itemize}

配置路径：\texttt{gateway.*}

\section{启动方式}

\subsection{前台运行}

\begin{lstlisting}[language=bash]
openclaw gateway --port 18789
\end{lstlisting}

\subsection{守护进程模式}

\begin{lstlisting}[language=bash]
openclaw daemon install   # Install systemd/launchd service
openclaw daemon start     # Start the daemon
openclaw daemon status    # Check status
\end{lstlisting}

\section{智能体循环（Agent Loop）}

每条入站消息触发一次完整的智能体循环：

\begin{center}
\begin{tikzpicture}[
    node distance=0.6cm,
    agentstep/.style={draw=teal!60, fill=teal!8, rounded corners=4pt, minimum width=12cm, minimum height=0.7cm, font=\small, align=left},
    arr/.style={-{Stealth[length=2.5mm]}, thick, draw=teal!60}
]
    \node[agentstep] (s1) {1. 接收入站消息，验证参数，解析会话};
    \node[agentstep, below=of s1] (s2) {2. 路由到目标 Agent，加载 Skills 快照};
    \node[agentstep, below=of s2] (s3) {3. 组装系统提示（基础提示 + Skills 提示 + 引导上下文）};
    \node[agentstep, below=of s3] (s4) {4. 调用模型推理，流式输出助手回复};
    \node[agentstep, below=of s4] (s5) {5. 执行工具调用（可能在沙箱中）};
    \node[agentstep, below=of s5] (s6) {6. 持久化会话，发送回复到渠道};

    \draw[arr] (s1) -- (s2);
    \draw[arr] (s2) -- (s3);
    \draw[arr] (s3) -- (s4);
    \draw[arr] (s4) -- (s5);
    \draw[arr] (s5) -- (s6);
\end{tikzpicture}
\end{center}

\section{RPC 接口}

Gateway 网关通过 WebSocket 和 HTTP 暴露 RPC 接口：

\begin{itemize}
  \item \texttt{agent} / \texttt{agent.wait}：触发智能体运行
  \item \texttt{sessions.list}：列出会话
  \item \texttt{voicecall.*}：语音通话（插件）
  \item \texttt{message.send}：发送消息到渠道
\end{itemize}

示例（通过 CLI 调用 RPC）：

\begin{lstlisting}[language=bash]
openclaw gateway call sessions.list --params '{}'
\end{lstlisting}

\section{认证}

推荐使用 \textbf{Anthropic API 密钥}：

\begin{lstlisting}[language=bash]
export ANTHROPIC_API_KEY="..."
# Or store in daemon env file:
echo 'ANTHROPIC_API_KEY=...' >> ~/.openclaw/.env
\end{lstlisting}

也支持 OAuth 和 setup-token 流程。使用 \texttt{openclaw models status} 验证认证状态。

\section{常用运维命令}

\begin{lstlisting}[language=bash]
openclaw doctor            # Health check
openclaw status            # Show storage paths & recent sessions
openclaw models status     # Check model auth
openclaw gateway --port N  # Start with custom port
\end{lstlisting}
